% Worked example, version C: a tug of war between the two nearest traditions.
\documentclass[border=3pt]{standalone}
% Shared preamble of the worked-example figures (case_a/b/c.tex); data from ../generated/case_data.tex
\usepackage{libertinus}
\renewcommand{\familydefault}{\sfdefault}
\usepackage{tikz}
\usepackage{amsmath}
\usepackage{booktabs}
\usetikzlibrary{arrows.meta,positioning,calc,fit,backgrounds,decorations.pathreplacing}
\makeatletter
\newcommand{\C}[1]{\@ifundefined{case@#1}{\textcolor{red}{[#1?]}}{\csname case@#1\endcsname}}
\makeatother
% word-class palette (decision_shares.png / playground), one slot per category
\definecolor{catF}{HTML}{2A78D6}  % function words
\definecolor{catV}{HTML}{EB6834}  % verbs
\definecolor{catN}{HTML}{1BAF7A}  % nouns and realia
\definecolor{catA}{HTML}{EDA100}  % agents and deities
\definecolor{catX}{HTML}{E87BA4}  % foreign words
\definecolor{catO}{HTML}{9A9A94}  % other / unclassified
\definecolor{shade}{HTML}{EB6834}
\definecolor{ink}{HTML}{2B2B2B}
\definecolor{soft}{HTML}{8A8A8A}
\definecolor{topc}{HTML}{1F7A6B}
\definecolor{secc}{HTML}{5B4B8A}
% a word tinted by its class; dashed frame when the class is low-confidence
\tikzset{wc-low/.style={densely dashed}, wc-ok/.style={}, reset/.style={text width={}, minimum height=0pt, minimum width=0pt, align=center, draw=none, fill=none, solid, font=}}
\newcommand{\wc}[3]{\tikz[baseline=(w.base)]\node[reset, inner sep=0.8pt, rounded corners=1pt, fill=#1!22,
  draw=#1!70, line width=0.4pt, wc-#2](w){\strut #3};}
% a word shaded by its share of the decision (0-100)
\newcommand{\ws}[2]{\tikz[baseline=(w.base)]\node[reset, inner sep=0.8pt, rounded corners=1pt, fill=shade!#1](w){\strut #2};}
\newcommand{\broken}[1]{\textcolor{soft}{#1}}
\newcommand{\tline}[2]{\noindent\hangindent=5.7mm\hangafter=1\makebox[4.5mm][r]{\textcolor{soft}{#1}}\hspace{1.2mm}#2\par}
\newcommand{\wf}[1]{\textbf{#1}}          % a word form
\newcommand{\sg}[1]{\textit{#1}}          % a sign sequence
\newcommand{\stepno}[1]{\tikz[baseline=-0.6ex]\node[reset, circle, fill=ink, text=white, inner sep=0.6pt, font=\scriptsize\bfseries]{#1};}
\newcommand{\legend}{%
  \wc{catF}{ok}{function} \wc{catV}{ok}{verb} \wc{catN}{ok}{noun, realia} \wc{catA}{ok}{agent, deity}
  \wc{catX}{ok}{foreign} \wc{catO}{low}{unclear}}
% generated by make_case_figures.py; do not edit
\expandafter\def\csname case@txtid\endcsname{KBo 41.22}
\expandafter\def\csname case@filed\endcsname{470}
\expandafter\def\csname case@words\endcsname{33}
\expandafter\def\csname case@top\endcsname{Luwian rituals (Puriyanni, Kuwattalla)}
\expandafter\def\csname case@topshort\endcsname{Luwian rituals}
\expandafter\def\csname case@second\endcsname{Kizzuwatna rituals}
\expandafter\def\csname case@secondshort\endcsname{Kizzuwatna rituals}
\expandafter\def\csname case@s1\endcsname{0.39}
\expandafter\def\csname case@s2\endcsname{0.25}
\expandafter\def\csname case@margin\endcsname{0.14}
\expandafter\def\csname case@tau\endcsname{0.018}
\expandafter\def\csname case@nfeat\endcsname{177}
\expandafter\def\csname case@dim\endcsname{25{,}372}
\expandafter\def\csname case@ntop\endcsname{10}
\expandafter\def\csname case@nsecond\endcsname{165}
\expandafter\def\csname case@common\endcsname{−21\%}
\expandafter\def\csname case@idf.rare\endcsname{6.8}
\expandafter\def\csname case@idf.common\endcsname{1.2}
\expandafter\def\csname case@nn1\endcsname{KBo 9.125+}
\expandafter\def\csname case@nn1cth\endcsname{409}
\expandafter\def\csname case@nn.409\endcsname{5}
\def\CaseText{\tline{3′}{\wc{catJ}{ok}{i-da-a-lu-un} \wc{catO}{low}{UD-an} \wc{catO}{low}{ma-}}\tline{4′}{\wc{catA}{ok}{DINGIR\textsuperscript{MEŠ}-aš} \wc{catT}{ok}{kar-pí-in} \wc{catO}{low}{pa-an-ga-u-wa-aš} \wc{catO}{low}{EME-an}}\tline{5′}{\wc{catV}{ok}{ša-me-ed-du-wa-ra-at-kán} \broken{ḫi-x}}\tline{6′}{\wc{catF}{ok}{na-aš-ta} \wc{catA}{ok}{EN} \wc{catN}{ok}{SÍSKUR} \wc{catT}{ok}{ŠUM-ŠU}}\tline{7′}{\wc{catF}{ok}{le-e} \wc{catO}{low}{ku-wa-pí-ik-ki} \broken{a-x}}\tline{8′}{\wc{catT}{ok}{te-eš-ḫa-an} \wc{catO}{low}{ma-ni-ku-wa-an(-)}}\tline{9′}{\wc{catF}{ok}{nu-uš-ši} \wc{catN}{ok}{\textsuperscript{GI}pát-tar} \wc{catF}{ok}{pa-ra-a}}\tline{10′}{\wc{catA}{ok}{\textsuperscript{MUNUS}ŠU.GI-ma} \wc{catV}{ok}{me-ma-i} \broken{x}}\tline{11′}{\wc{catA}{ok}{DINGIR\textsuperscript{MEŠ}-aš} \wc{catT}{ok}{kar-pí-in} \wc{catO}{low}{pa-an-ga-u-wa-aš}}\tline{12′}{\wc{catA}{ok}{\textsuperscript{LÚ}A.ZU-ma-aš-ši-iš-ša-an}}\tline{13′}{\wc{catF}{ok}{na-aš} \wc{catF}{ok}{ša-ra-a} \wc{catO}{low}{ti-}}\tline{14′}{\wc{catF}{ok}{na-aš-ta} \wc{catB}{ok}{GU₄-un}}\tline{15′}{\wc{catF}{ok}{a-ra-aḫ-za-an-da}}}
\def\CaseShaded{\tline{3′}{\ws{27}{i-da-a-lu-un} \ws{2}{UD-an} \ws{0}{ma-}}\tline{4′}{\ws{2}{DINGIR\textsuperscript{MEŠ}-aš} \ws{39}{kar-pí-in} \ws{70}{pa-an-ga-u-wa-aš} \ws{93}{EME-an}}\tline{5′}{\ws{0}{ša-me-ed-du-wa-ra-at-kán} \broken{ḫi-x}}\tline{6′}{\ws{0}{na-aš-ta} \ws{0}{EN} \ws{0}{SÍSKUR} \ws{14}{ŠUM-ŠU}}\tline{7′}{\ws{0}{le-e} \ws{15}{ku-wa-pí-ik-ki} \broken{a-x}}\tline{8′}{\ws{0}{te-eš-ḫa-an} \ws{6}{ma-ni-ku-wa-an(-)}}\tline{9′}{\ws{0}{nu-uš-ši} \ws{0}{\textsuperscript{GI}pát-tar} \ws{2}{pa-ra-a}}\tline{10′}{\ws{27}{\textsuperscript{MUNUS}ŠU.GI-ma} \ws{0}{me-ma-i} \broken{x}}\tline{11′}{\ws{2}{DINGIR\textsuperscript{MEŠ}-aš} \ws{39}{kar-pí-in} \ws{70}{pa-an-ga-u-wa-aš}}\tline{12′}{\ws{0}{\textsuperscript{LÚ}A.ZU-ma-aš-ši-iš-ša-an}}\tline{13′}{\ws{0}{na-aš} \ws{0}{ša-ra-a} \ws{5}{ti-}}\tline{14′}{\ws{0}{na-aš-ta} \ws{0}{GU₄-un}}\tline{15′}{\ws{0}{a-ra-aḫ-za-an-da}}}
\def\CaseKeyWord{\wc{catO}{low}{pa-an-ga-u-wa-aš}}
\def\CaseBars{\simbar{0}{Luwian rituals}{0.389}{top}\simbar{1}{Kizzuwatna rituals}{0.245}{second}\simbar{2}{Hattian cult layer}{0.204}{other}\simbar{3}{Syro-Mesopotamia}{0.181}{other}\simbar{4}{Kizzuwatna–Lower Land}{0.173}{other}\simbar{5}{Palaic festivals}{0.165}{other}\simbar{6}{Luwian festivals}{0.165}{other}\simbar{7}{Ḫuwaššanna festivals}{0.160}{other}\simbar{8}{Mixed Hurrian}{0.147}{other}\simbar{9}{Arzawa}{0.146}{other}}
\def\CaseFeatRows{\wf{EMEan} & 0.11 & 0.17 & 0.00 & +18.1 \\\sg{EME-an} & 0.10 & 0.17 & 0.00 & +17.1 \\\wf{pankauaš} & 0.18 & 0.08 & 0.00 & +14.7 \\\sg{ga-u-wa} & 0.17 & 0.08 & 0.00 & +12.7 \\\sg{an-ga-u} & 0.16 & 0.07 & 0.00 & +11.8 \\\sg{pa-an-ga} & 0.16 & 0.07 & 0.00 & +11.8 \\\midrule \wf{EN} & 0.05 & 0.00 & 0.06 & −3.0 \\\sg{EN} & 0.04 & 0.00 & 0.06 & −2.3 \\}
\def\CaseTug{\tug{0}{\wf{EMEan}}{18.1}\tug{1}{\sg{EME-an}}{17.1}\tug{2}{\wf{pankauaš}}{14.7}\tug{3}{\sg{ga-u-wa}}{12.7}\tug{4}{\sg{an-ga-u}}{11.8}\tug{5}{\sg{pa-an-ga}}{11.8}\tug{6}{\wf{EN}}{-3.0}\tug{7}{\sg{EN}}{-2.3}}
\expandafter\def\csname case@x.w\endcsname{0.18}
\expandafter\def\csname case@x.n1\endcsname{0.16}
\expandafter\def\csname case@x.n2\endcsname{0.16}
\expandafter\def\csname case@x.n3\endcsname{0.17}


\begin{document}
\begin{tikzpicture}[font=\sffamily\scriptsize\color{ink},
  box/.style={draw=soft!60, rounded corners=3pt, line width=0.5pt, inner sep=4pt, fill=white},
  title/.style={font=\sffamily\footnotesize\bfseries, anchor=south west, inner sep=1pt}]
% the fragment, centred at the top
\node[box, text width=72mm, anchor=north] (frag) at (0,0) {%
  {\linespread{1.12}\selectfont\raggedright\CaseText}\vspace{1pt}{\tiny\legend}};
\node[title] at (frag.north west) {\C{txtid}: an orphan (CTH \C{filed}), \C{words} legible words};
% the two traditions as poles
\node[box, draw=topc, fill=topc!8, text width=38mm, align=center, anchor=north east] (L) at (-44mm,-3mm) {%
  \textbf{\color{topc}\C{top}}\\ \C{ntop} reference manuscripts\\[2pt] $s_1=\cos(\mathbf{x},\mathbf{c}_1)=\textbf{\C{s1}}$};
\node[box, draw=secc, fill=secc!8, text width=38mm, align=center, anchor=north west] (R) at (44mm,-3mm) {%
  \textbf{\color{secc}\C{second}}\\ \C{nsecond} reference manuscripts\\[2pt] $s_2=\cos(\mathbf{x},\mathbf{c}_2)=\textbf{\C{s2}}$};
\node[font=\tiny, text width=38mm, align=center, anchor=north] at (L.south) {nearest tradition};
\node[font=\tiny, text width=38mm, align=center, anchor=north] at (R.south) {runner-up};
% the tug: each feature pulls toward the tradition where it is more typical
\begin{scope}[shift={($(frag.south)+(0,-9mm)$)}, x=3.2mm, y=-4.2mm]
  \node[title, anchor=south] at (0,-0.6) {Which features pull where: $x_f\,(c_{1,f}-c_{2,f})\times10^3$};
  \newcommand{\tug}[3]{%
    \pgfmathsetmacro{\v}{#3}
    \pgfmathparse{\v>0 ? 1 : 0}
    \ifnum\pgfmathresult=1
      \fill[topc!80] (0,#1+0.15) rectangle (-\v,#1+0.85);
      \node[anchor=west, inner sep=1pt] at (0.2,#1+0.5) {#2};
      \node[anchor=east, inner sep=1pt, font=\tiny] at (-\v-0.1,#1+0.5) {\pgfmathprintnumber[fixed,precision=1]{\v}};
    \else
      \fill[secc!80] (0,#1+0.15) rectangle (-\v,#1+0.85);
      \node[anchor=east, inner sep=1pt] at (-0.2,#1+0.5) {#2};
      \node[anchor=west, inner sep=1pt, font=\tiny] at (-\v+0.1,#1+0.5) {\pgfmathprintnumber[fixed,precision=1]{\v}};
    \fi}
  \CaseTug
  \draw[ink, line width=0.6pt] (0,-0.1) -- (0,8.2);
  \node[anchor=north east, text=topc, font=\tiny\bfseries] at (-0.3,8.3) {$\leftarrow$ toward \C{topshort}};
  \node[anchor=north west, text=secc, font=\tiny\bfseries] at (0.3,8.3) {toward \C{secondshort} $\rightarrow$};
\end{scope}
% decision and check
\node[box, text width=60mm, anchor=north, align=center] (dec) at ($(frag.south)+(0,-50mm)$) {%
  All pulls add up to the lead of the nearest tradition:\\[2pt] $m=\sum_f x_f(c_{1,f}-c_{2,f})=s_1-s_2=\textbf{\C{margin}}$\\[2pt]
  $m \geq \tau=\C{tau}$ $\Rightarrow$ placed with \textbf{\color{topc}\C{topshort}}\\
  {\tiny ($\tau$: the lead at which 80\% of test fragments are placed correctly)}};
\node[box, text width=40mm, anchor=north west, draw=soft] (how) at ($(L.south west)+(0,-6mm)$) {%
  \textbf{How a fragment becomes numbers.} Every word form and every sign sequence in a word is a feature $f$, weighted $x_f\propto(1+\log\mathit{tf}_f)\,\mathit{idf}_f$: frequent here, rare in the corpus. \wf{pankauaš}: idf \C{idf.rare}; the sign \sg{na}: \C{idf.common}. Each tradition is the average $\mathbf{c}_t$ of its manuscripts.};
\node[box, text width=40mm, anchor=north east, draw=soft] (chk) at ($(R.south east)+(0,-6mm)$) {%
  \textbf{Check by hand.} Of the five single manuscripts nearest to it, \C{nn.409} are Tunnawiya's ritual (CTH \C{nn1cth}), which no reference group includes: a lead for the philologist.};
\end{tikzpicture}
\end{document}
